\documentclass[12pt,a4paper,oneside]{article}
\usepackage[brazil]{babel}
\usepackage[utf8]{inputenc}
\usepackage{olymp}
\usepackage{color}
\usepackage[dvipsnames]{xcolor}
\usepackage{listings}

\setcounter{problem}{1}

\begin{document}

\begin{problem}{Divisão inteira e seu resto (versão 2)}{divresto}{2}

Faça uma função que calcule a divisão inteira entre dois números naturais e também seu resto.

A função deve obrigatoriamente ter o seguinte protótipo:

%\texttt{\textcolor{YellowGreen}{\textbackslash\textbackslash xxx}}\\
\texttt{\textcolor{Mulberry}{bool} divresto(\textcolor{Mulberry}{int}, \textcolor{Mulberry}{int}, \textcolor{Mulberry}{int \&}, \textcolor{Mulberry}{int \&})};

Os primeiros dois parâmetros são os valores do dividendo e do divisor, recebidos pela função. Os dois últimos são os valores da divisão inteira e do resto da divisão, retornados por referência. A função deve ainda retornar, por valor, verdadeiro se a divisão for definida ou falso caso contrário.

Faça um programa que usa essa função para realizar a operação em vários casos de teste.

%\Notes
%
%Observações, se tiver.

\InputFile

A entrada contém múltiplos casos de teste. A primeira linha contém um número $N$, indicando o número de testes que serão realizados. As próximas $N$ linhas contêm pares de valores naturais $A$ e $B$, respectivamente o dividendo e o divisor da operação. Restrições: $1 \leq N \leq 1000$, $0 \leq A,B \leq 100000$.

\OutputFile

Seu programa deve gerar $N$ linhas na saída, uma para cada par de valores lidos $A$ e $B$ na entrada, contendo o quociente (inteiro) da divisão de $A$ por $B$, um espaço em branco e o resto da divisão inteira de $A$ por $B$. Caso a operação não esteja definida, escreva ``\texttt{indefinido}''.


 \Notes
Você pode escrever os resultados à medida que lê a entrada, não precisa ler a entrada toda antes.

\Example

\begin{example}
\exmp{
5
25 5
29 3
5 10
100 0
342 674
}{
5 0
9 2
0 5
indefinido
0 342
}%
\end{example}


%Comentários sobre o exemplo, se necessário.

\end{problem}

\end{document}
